\subsection{A compulsory integral obstruction and the exact local calculation at two}

\begin{statement}{B.35}{Every good prime has a nonzero determinant-extension class}{Proved}
For this statement and B.28, B.31--B.32 use the tangent at
$O=C_0=[1:-1:0]$ with unit parameter $q_0=Z/X$, exactly as
in B.17 and B.24. This is separate from the $A_0$ origin
of B.29--B.30. Let $\eta$ be the resulting integral class
of the extension of the primitive compact determinant
$w$, and let $m_p$ be its local order exponent (B.31.3).
Then
\[
                         m_p\ge1\qquad(p\ne3).                \tag{B.35.1}
\]
In particular, for every $p\ne3$ the actual zero-abelian
class-two section $\kappa_3(p)w$ does not lift to class three.
Every liftable input of (B.31.1) must have $3\mid a_p$ at
every prime away from three.

At two the actual local class is nonzero even after
quotienting by the image of
\[
 \operatorname{ad}_w:H^1_{\rm ur}(K_{(2)},V)
           \longrightarrow H^1_{\rm ur}(K_{(2)},A(-1)),
       \qquad v\longmapsto[v,w].                              \tag{B.35.2}
\]
This is an integral assertion; it does not prove that the
rational Ross edge is a nonzero arithmetic generator.

There is an exact smaller calculation for $m_2$.
Quaternion automorphisms reduce its coefficient lattice
from rank $330$ to rank $41$.
Over $\mathbb F_4$ the resulting cokernel is
\[
 C_4=(\mathbb Z/3)^9\oplus(\mathbb Z/9)^4
       \oplus(\mathbb Z/27)^8\oplus(\mathbb Z/81)^4.
                                                               \tag{B.35.3}
\]
Keeping the actual Frobenius-two descent gives instead
\[
 C_2=(\mathbb Z/3)^5\oplus(\mathbb Z/9)^2
       \oplus(\mathbb Z/27)^4\oplus(\mathbb Z/81)^2.
                                                               \tag{B.35.4}
\]
The tangent-rescaling equation and that descent restrict
$\eta$ to an explicitly specified affine set of $3^{10}$
classes. Its difference group is
\[
                (\mathbb Z/3)^4\oplus\mathbb Z/9
                           \oplus\mathbb Z/81.                \tag{B.35.5}
\]
This affine set contains module classes of each order
$3,9,27,81$, certified below. Consequently these exact
symmetry conditions alone do not determine which of
$1,2,3,4$ is the actual $m_2$.
B.36 adds the actual tangential translation data and
settles this value as $m_2=3$.
No member of this finite affine set is asserted to be a
geometric action merely because it satisfies the linear conditions.
\end{statement}
\begin{proof}
\emph{The integral tangent-rescaling identity.}
Work first over $K=\mathbb Q(\zeta_3)$ and use the genuine
curve automorphism
\[
              \sigma[X:Y:Z]=[X:Y:\zeta_3Z].                  \tag{B.35.6}
\]
It fixes $O$, preserves $w$, and scales its unit tangent
by $\zeta_3$.
Choose a geometric path identifying the image tangent
with the original one. The resulting based lift
$\widetilde\sigma$ and the genuine arithmetic action obey
\[
 \widetilde\sigma\rho_g\widetilde\sigma^{-1}
                  =\operatorname{Int}(k_\sigma(g))\rho_g,
 \qquad [k_\sigma^{\rm ab}]
                  =\kappa_3(\zeta_3)e_O .                    \tag{B.35.7}
\]
The latter equality is functoriality of the local
parameter's Kummer character on two tangents of ratio
$\zeta_3$.
If $\omega\in L=\gamma_2N/\gamma_4N$ lifts $w$, write
$\widetilde\sigma\omega=\omega+t_\sigma$, with
$t_\sigma\in A$.
Apply (B.35.7) to $\omega$ and use
$g\omega=\chi_3(g)(\omega+\eta(g))$.
After dividing by $\chi_3(g)$, the $t_\sigma$ term is
exactly a coboundary in $A(-1)$. Thus, with the positive
peripheral convention of B.17,
\[
       (\sigma-I)[\eta]
                   =\operatorname{ad}_{w*}
                         (\kappa_3(\zeta_3)e_O).             \tag{B.35.8}
\]
The action of $\sigma$ on $A$ has order three.
The chosen based lift need not have order three:
its cube may be peripheral inner. Such an inner change
is trivial on $A$ and does not alter (B.35.8).

Two primitive inclusions detect its right side.
First $\mathbb Z_3(1)e_O\subset V$ has free quotient,
as follows from the homology sequence with the toric
lattice $\mathbb Z_3(1)[D]/\mathbb Z_3(1)\sum e_d$.
Second, $\operatorname{ad}_w:V\to A(-1)$ is a primitive
injection. Indeed $w=u\wedge v$ is primitive decomposable,
so $u,v$ extend to a $\mathbb Z_3$-basis of $V$.
In an order beginning with these two vectors, the ten
$[e_i,w]$ are ten distinct integral Hall basis vectors.
Its quotient is therefore free.
At every good prime, the rational eigenvalues on both
quotients have complex absolute value strictly greater
than one: the weights of $V$ are one and two, and those
of $A(-1)$ are one through four.
Their invariant groups are zero.
The two cohomology exact sequences consequently give
injections
\[
 H^1_{\rm ur}(K_{\mathfrak q},\mathbb Z_3(1)e_O)
  \lhook\joinrel\longrightarrow H^1_{\rm ur}(K_{\mathfrak q},V)
  \lhook\joinrel\longrightarrow H^1_{\rm ur}(K_{\mathfrak q},A(-1)).
                                                               \tag{B.35.9}
\]
For $Q=N\mathfrak q$, the first group is
$\mathbb Z_3/(Q-1)$, and $\kappa_3(\zeta_3)$ has exact
order three in it. Equivalently its Frobenius value is
$(Q-1)/3$ times a three-adic unit, modulo $Q-1$.
This remains true when its first few finite Kummer
reductions vanish because $\zeta_3$ happens to have
three-power roots in the residue field.
The full completed class still has order three.
Thus the right side of (B.35.8) is nonzero locally at
every $\mathfrak q\nmid3$. This proves (B.35.1), and
B.31 gives the stated lifting consequences.

\emph{The actual characteristic-two curve and quaternion action.}
Reduce the smooth marked curve at two and pass to
$\mathbb F_4=\mathbb F_2[\zeta]/(\zeta^2+\zeta+1)$.
The projective change of coordinates
\[
             x=\frac Z{X+Y},\qquad
             y=\frac X{X+Y}+\zeta
                                                               \tag{B.35.10}
\]
gives $E_0:y^2+y=x^3$, with $O$ at infinity.
Its eight other $\mathbb F_4$ points are precisely the
other coordinate cusps. In the order of (B.24.1) they are
\[
 (1,\zeta),\ (\zeta^2,\zeta),\ (\zeta,\zeta),\
 (1,\zeta^2),\ (\zeta^2,\zeta^2),\ (\zeta,\zeta^2),\
 (0,0),\ (0,1).
                                                               \tag{B.35.11}
\]
The tangent-line addition formula gives
$[2](x,y)=(x^4,y^4+1)$. Hence Frobenius four is $[-2]$.
All nine points are $E_0[3]$, and their coordinates in
basis $A_0,A_2$ are exactly those used to construct $U$
in B.24. No comparison from a different cusp lattice is used.

There are eight automorphisms
\[
 a_{r,b}(x,y)=(x+r,y+r^2x+b),\qquad
 r\in\mathbb F_4,\quad b^2+b=r^3.                            \tag{B.35.12}
\]
Substitution verifies the equation. For $r\ne0$ their
square is $(x,y)\mapsto(x,y+1)$; for $r=0$ they are
the identity and that involution. Their composition law
therefore gives the quaternion group $Q_8$.
They fix $O$ and have derivative one in $t=x/y$.
Since $t=q_0(1+O(q_0^3))$, this is the same unit tangent
as before reduction.
The induced tangential actions are genuine: a substitution
$t\mapsto t\,u(t)$ with $u(0)=1$ has unique compatible
three-power roots of $u(t)$, so its maps on the prime-to-two
Puiseux fibres compose literally.
These automorphisms are defined over $\mathbb F_4$ and
commute with its based arithmetic Frobenius.
Prime-to-two specialization transports the actual action
and tangent from the original local curve
\citep[XIII, \S\S2.9--2.12]{BSGA1}.

Average $\omega$ over $Q_8$ with $1/8$. This is integral,
still lifts $w$, and makes its Frobenius-four difference
lie in $A(-1)^{Q_8}$. The averaging projector commutes
with Frobenius, so its cokernel embeds as a direct summand
of the original local cokernel; the order of the actual
class is preserved.
On compact homology the quaternion group has its
rank-two faithful representation. On the eight-dimensional
boundary lattice it is its regular permutation representation.
Thus the traces on $V$ are $10$ at one, $-2$ at the central
involution, and zero at the six order-four elements.
The degree-three Witt trace gives
\[
                    \operatorname{rank}A^{Q_8}
                              =(330-2)/8=41.                  \tag{B.35.13}
\]

\emph{All matrices and the Frobenius-two descent.}
Let $s\in\mathbb Z_3$ be the root of $s^2=-2$ congruent
to one modulo three, and put
\[
 I=\begin{pmatrix}s&-1\\-1&-s\end{pmatrix},\qquad
 J=\begin{pmatrix}1&s\\s&-1\end{pmatrix},\qquad K=IJ.
                                                               \tag{B.35.14}
\]
These are compact quaternion generators.
For any one of their eight products, denote its matrix
by $M$ and its actual cusp permutation by $P_M$.
The corresponding integral degree-one matrix is
\[
                  Q_M=\begin{pmatrix}
                   M&0\\(UM-P_MU)/3&P_M
                  \end{pmatrix}.                             \tag{B.35.15}
\]
The residue action on the displayed cusp basis makes the
numerator divisible by three.
This realizes the actual representation up to a strict
integral change of basis. To prove that assertion, lift
a conjugacy of the two residual quaternion representations
through the successive three-primary congruence kernels.
Their first cohomology for $Q_8$ vanishes by averaging
with $1/8$, so the conjugacies lift compatibly.
Multiplying the resulting conjugator by the unique scalar
in $1+3\mathbb Z_3$ whose square is its inverse determinant
keeps that conjugacy and makes its determinant one.
Thus the specified vector $w$ is unchanged as well.
For the resulting compact change $B\equiv I\pmod3$, the
full change is
\[
                   \begin{pmatrix}B&0\\U(B-I)/3&I_8\end{pmatrix}.
\]
It is integral and commutes with the Frobenius-four matrix
$\left(\begin{smallmatrix}-2I_2&0\\-2U&4I_8\end{smallmatrix}\right)$.
Thus the ensuing cokernel computation is independent of
that strict conjugacy.

The compatible compact matrix for (B.35.6) is
\[
                    S=-\tfrac12(I_2+I+J+K).                  \tag{B.35.16}
\]
It has order three, normalizes $Q_8$, and reduces to $R_3$.
After the preceding strict conjugacy the actual matrix
must be this one: both induce the same conjugation on
$Q_8$, their ratio is scalar, and $\mathbb Q_3$ contains
no nontrivial scalar cube root of unity.
On the original $\mathbb F_2$ curve, Frobenius two becomes
$(x,y)\mapsto(x^2,y^2+1)$ in (B.35.10), including the
extra one from the $\zeta$-dependent coordinate change.
Its compatible compact matrix is $J-K$.
It has square $-2I$, determinant two, and residual matrix
$C_E$ of (B.28.4). The quaternion conjugation and determinant
identify it uniquely; the possible scalar sign is fixed
by that residual matrix.
Its full matrix is
\[
 F_2=\begin{pmatrix}
 J-K&0\\(U(J-K)-2P_cU)/3&2P_c
 \end{pmatrix},\qquad
 F_2^2=\begin{pmatrix}-2I&0\\-2U&4I\end{pmatrix}.
                                                               \tag{B.35.17}
\]
Here $P_c$ exchanges the subscripts one and two.

Take the third Lie functors and apply the integral
quaternion projector. On its rank-$41$ image write
\[
 R_2=\tfrac12\operatorname{Lie}_3(F_2),\quad
 D_4=R_2^2-I,\quad D_2=R_2-I,
 \quad S_A=\operatorname{Lie}_3(Q_S).                          \tag{B.35.18}
\]
The fresh program \texttt{b\_char2\_symmetry.py} constructs
these matrices from the finite-field automorphisms and
$U$, and gives the exact Smith valuations
\[
 \begin{array}{c|rrrrr}
       &0&1&2&3&4\\ \hline
 D_4&16&9&4&8&4\\
 D_2&28&5&2&4&2
 \end{array}.                                                  \tag{B.35.19}
\]
All pivots are below the working modulus $3^{10}$.
The quaternion projector has rank $41$ already modulo
three, and the rational positive weights give exact
full rank of both $D_i$, so the certificates determine
the complete integral cokernels. For $D_4$ the independent
multidegree determinant valuation is $4+2\cdot16+21=57$.
This proves (B.35.3)--(B.35.4).
Restriction from the unramified quadratic extension
sends a class represented by $x$ to $(I+R_2)x$.
Its image is $C_4^{R_2}$ and its kernel is zero, because
the restriction/corestriction composite is two.
It preserves the order of the local extension class.

\emph{The exact residual finite ambiguity.}
Set $b=[e_O,w]$ in the rank-$41$ lattice; it is fixed
by $Q_8$. The nonzero class of $b$ in $C_4$ has order
three, by (B.35.9). Equation (B.35.8) is the finite
inhomogeneous system
\[
        (S_A-I)(I+R_2)x=b\quad\pmod{D_4A^{Q_8}},
             \qquad x\in C_2.                               \tag{B.35.20}
\]
The same program computes its homogeneous group exactly.
For example, the dual of the simultaneous fixed subgroup
is presented by the concatenated matrix
\[
                [\,D_4^{\rm t},\ S_A^{\rm t}-I,
                                      \ R_2^{\rm t}-I\,].       \tag{B.35.21}
\]
The transpose-inverse differences have the same images
as the displayed transpose differences, and all actions
commute with $D_4$. Its nonzero Smith valuations are
$1^4,2^1,4^1$, with $35$ unit pivots. This proves the
finite difference group (B.35.5).

In the positive Smith coordinates of $C_2$, whose
exponents are
$(1,1,1,1,1,2,2,3,3,3,3,4,4)$,
the following four vectors solve (B.35.20):
\[
 \begin{array}{c|l}
 \text{order}&\text{coordinates}\\ \hline
 3&(0,0,2,1,2,3,0,0,0,0,18,0,0)\\
 9&(0,0,2,1,2,3,0,9,18,18,0,45,27)\\
 27&(0,0,2,1,2,3,0,3,15,6,3,15,9)\\
 81&(0,0,0,2,2,3,0,19,14,20,4,59,30)
 \end{array}.                                                  \tag{B.35.22}
\]
The exact Smith-coordinate algorithm, all transport
matrices modulo $3^{10}$, and the machine-readable vectors
are retained in the companion code and data.
The program verifies both the equations and the four
exact orders. They are finite module classes, not a
claim of four geometric curves or four possible actual
actions of this one curve. They prove the stated
limitation of the tested conditions.

Finally suppose the actual local class were in the
image of (B.35.2). That map is injective, and its image
map is equivariant. Quaternion invariance would put its
unique antecedent in $H^1_{\rm ur}(K_{(2)},V)^{Q_8}$.
Exact averaging and the regular permutation description
give
\[
 V^{Q_8}=\mathbb Z_3(1)e_O,\qquad
 H^1_{\rm ur}(K_{(2)},V)^{Q_8}
                        =(\mathbb Z/3)e_O.
\]
The automorphism $\sigma$ fixes this subgroup, whereas
(B.35.8) has its nonzero right side. This is a
contradiction. Thus the actual integral class remains
nonzero modulo the inner/path image, proving the
additional assertion of the statement. Selecting its
particular member of (B.35.20), rather than enlarging
the finite test without an arithmetic comparison, is
the remaining local problem.
\end{proof}
